\section{Intuitionistic Proof Transformation}

\subsection{Induction principle for path relation}

Define
\[
G_jx \equiv \forall \Vec z(
((x)_0,\ldots,(x)_{l-1}) = \Vec v \imp
(\Gamma_{2j})^\circ_{\Vec v} \imp \Delta_{2j}^\bullet), \\
Gx_0x \equiv \Land_{j \in K} (x_0=j \imp G_jx),
\]
where 
$l$ is the number of inductive atomic formulas in $\Gamma_{2j}$, and
$\Vec z$ is $\FV((\Gamma_{2j})^\circ_{\Vec v},\Delta_{2j}^\bullet)$.

$>_\Pi$ is defined in the appendix.
%
We define $\rho$ as obtained in Theorem \ref{th:rankingfunc}.

We define $\rho'(c,x)$ as
$\rho(c,((x)_0,\ldots,(x)_{l-1}))$
where $l$ is the number of inductive atomic formulas in 
the antecedent of the sequent of companion number $c$.

\begin{Prop}\label{prop:0.5}
For a path $\pi$ from a companion $c$ to a bud $d$ above $c$,
$(c,x) >_\pi (d,y)$ implies $\rho'(c,x) >_\Lex \rho'(d,y)$.
\end{Prop}

{\em Proof.}
In a given cyclic proof $\Pi$,
for a companion $c_i$ and a bud $d_i$ above $c_i$,
we write a path $\pi_i$ for the path from $c_i$ to $d_i$.

For $\pi_i$, we define an SC graph $G_{\pi_i}$ 
as $(c_i,(S_i,P_i),d_i)$ where
$j S_i l$ is defined to hold if there is some nonprogressing trace from the $j$-th
inductive atomic formula to the $l$-th inductive atomic formula, and
$j P_i l$ is defined to hold if there is some progressing trace from the $j$-th
inductive atomic formula to the $l$-th inductive atomic formula.

The set of SC graphs $G_{\pi_i}$ satisfies the SCT condition.
It is shown as follows.
Assume an infinite composable sequence $G_{\pi_{p_1}},G_{\pi_{p_2}},\ldots$
of composable SC graphs.
Then $pi_{p_1},\pi_{p_2},\ldots$ is an infinite path in the cyclic proof.
By the global trace condition, there is some trace for some tail of the path
that is infinitely progressing.
Let $A_1,A_2,\ldots$ be the occurrences of the inductive atomic formulas
in the trace.
Take $n_i$ such that
$A_i$ is the $n_i$-th inductive atomic formula in the antecedent.
Then $n_1,n_2,\ldots$ is an SC trace of some tail of $G_{\pi_{p_1}},G_{\pi_{p_2}},\ldots$ that is infinitely progressing.
Hence the set of SC graphs satisfies the SCT condition.

By Theorem \ref{th:rankingfunc}, there is some $\rho$ such that
$(c,\Vec t) \to^{G_i} (d,\Vec u)$ implies
$\rho(c,\Vec t) >_\Lex \rho(d,\Vec u)$.

Fix  path $\pi$ from a companion $c$ to a bud $d$ above $c$.
Assume $(c,x) >_\pi (d,y)$.
Then $(x)_i > (y)_j$ iff there is a progressing trace from $i$ to $j$
iff $i P j$.
Similarly $(x)_i \ge (y)_j$ iff there is a nonprogressing trace from $i$ to $j$
iff $i S j$.
Hence we have
\[
(c,((x)_0,\ldots,(x)_{l-1})) \to^{G_\pi} (d,((y)_0,\ldots,(y)_{l'-1})).
\]

By the ranking function $\rho$, we have
$\rho(c,((x)_0,\ldots,(x)_{l-1})) >_\Lex \rho(d,((y)_0,\ldots,(y)_{l'-1}))$.
Hence $\rho'(c,x) >_\Lex \rho'(d,y)$.
$\Box$

\begin{Prop}\label{prop:1}
$(c,x) >_\Pi (d,y)$ implies $\rho'(c,x) >_\Lex \rho'(d,y)$.
\end{Prop}

The proof is given in the appendix.

\begin{Prop}\label{prop:2}
\[
\HA \prove 
\forall cx(\forall dy(\rho'(d,y) <_U \rho'(c,x) \imp Gdy) \imp Gcx)
\imp \forall cx.Gcx.
\]
\end{Prop}

The proof is given in the appendix

\begin{Prop}\label{prop:ind}
$\HA \prove \Ind(>_\pi,G)$.
\end{Prop}

{\em Proof.}
In the statement of Proposition \ref{prop:2},
since $\rho'(d,y) <_U \rho'(c,x)$ appears negatively,
by Proposition \ref{prop:1}, we can replace it by $(d,y) <_\Pi (c,x)$.
Hence $\Ind(>_\pi,G)$ is provable.
$\Box$

\subsection{Equivalence between cyclic proofs and inductive proofs}

Combining Proposition \ref{prop:ind} and
the proof transformation given in \cite{Berardi18}
we obtain the following theorem.
The detailed descriptions are given in the appendix.

\begin{Th}[Equivalence of $\LJIDHA$ and $\CLJIDHA$]
If $\CLJIDHA$ with inductive predicates $P_1,\ldots,P_n$ proves a sequent,
then there is a proof of the same sequent in $\LJIDHA$ with $P_1,\ldots,P_n$ and their stage-number inductive predicates $P_1',\ldots,P_n'$.
Moreover, we have an algorithm of
proof transformation from $\CLJIDHA$ to $\LJIDHA$.
Moreover, the correctness of the algorithm of the proof transformation
from $\CLJIDHA$ to $\LJIDHA$ is proved in intuitionistic logic.
\end{Th}
